\chapter{可预测性从哪里来}\label{ch:source}

\begin{keybox}[本章要回答的问题]
股票价格里到底有没有可被系统利用的规律？如果有，它来自哪里、能持续多久、
为什么没有被套利干净？以及——机器学习在其中究竟贡献了什么？
\end{keybox}

\section{一个思想实验：随机也能做出漂亮的回测}

在讲任何原理之前，先做一个实验。假设我们对 500 只股票的未来收益完全没有任何信息，
但做 1000 次独立的"策略研发"：每次随机给股票打分、按分数排序、取前 20\% 做多、
后 20\% 做空，回测五年。

结果会令你吃惊：这 1000 条净值曲线里，最好的一条年化收益可以轻松超过 20\%，
夏普比率超过 1.5。而且它看起来"很稳定"——月度胜率 60\% 以上，最大回撤不到 10\%。

这不是数据挖掘的偶然，而是极值分布的必然。若 1000 条纯噪声策略的夏普比率近似服从
$\N(0, \sigma^2)$，那么最大值的期望约为

\begin{equation}
  \E\!\left[\max_{1\le n\le N} \SR_n\right]
  \;\approx\; \sigma\left[(1-\gamma)\Phi^{-1}\!\left(1-\tfrac1N\right)
  + \gamma\,\Phi^{-1}\!\left(1-\tfrac{1}{Ne}\right)\right],
  \qquad \gamma \approx 0.5772,
  \label{eq:emax}
\end{equation}

其中 $\Phi^{-1}$ 是标准正态分位数。$N=1000$、$\sigma=1/\sqrt{T}$（$T$ 为样本天数）时，
即便真实夏普为 $0$，观测到的最大值也会显著大于零。\textbf{你看到的"最好策略"，
本质上是一个极值统计量，而不是一个能力证明。}

\begin{warnbox}[本书的第一个纪律]
任何回测结果，在知道"它是在多少次尝试中被选出来的"之前，都不构成证据。
第~\ref{ch:overfit} 章会把这个直觉变成可计算的检验（Deflated Sharpe Ratio 与 PBO）。
\end{warnbox}

\section{可预测性的五类真实来源}

抛开叙事，股价的可预测性只有有限的几个来源。理解它们，你才能判断一个"新因子"
是新发现，还是老故事的换皮。

\subsection{风险补偿：你赚的是承担风险的报酬}

有效市场理论并不否认"超额收益"，它否认的是"无风险的超额收益"。价值股长期跑赢，
可能只是因为它们在经济下行时跌得更惨，投资者要求补偿。这条逻辑链上有几类经典异象：

\begin{itemize}
  \item \textbf{价值}：账面市值比（BP）、盈利收益率（EP）高的股票，长期收益更高。
  \item \textbf{低波动}：低波动股票的风险调整后收益反而更高（杠杆约束与彩票偏好共同导致）。
  \item \textbf{质量}：高 ROE、低应计、现金流扎实的公司更抗跌。
  \item \textbf{流动性}：流动性差的股票需要更高的预期收益来补偿变现成本。
\end{itemize}

\textbf{关键点}：风险补偿型因子的收益是"周期性"的。它们在危机中会大幅回撤，
而且当越来越多资金涌向同一类股票时，补偿会被压缩——这就是"因子拥挤"。

\subsection{行为偏差：你赚的是别人犯错的代价}

A 股散户成交占比高，行为偏差的痕迹比成熟市场更明显：

\begin{itemize}
  \item \textbf{过度反应与反转}：短期（1--20 日）涨幅过大的股票，后续倾向回落。
        这是 A 股最强的横截面效应之一，本书实测的 5 日反转 Rank IC 在 $0.05$ 量级。
  \item \textbf{处置效应}：投资者倾向卖出盈利股、持有亏损股，造成动量在中长期偏弱。
  \item \textbf{彩票偏好}：正偏度（偶尔暴涨）的股票被系统性高估，
        因此"最大单日收益"是一个负向因子。
  \item \textbf{有限注意}：换手率骤升、登上龙虎榜的股票，短期被过度关注，随后回落。
\end{itemize}

\subsection{微观结构：交易规则本身制造了规律}

这一类最容易被忽视，却最"可执行"，因为它直接来自规则：

\begin{itemize}
  \item \textbf{T+1}：当日买入不能当日卖出，日内动量策略在 A 股基本不可行，
        但隔夜与次日开盘的行为因此更有规律。
  \item \textbf{涨跌停}：涨停日买不进、跌停日卖不出。任何假设"信号出现即可成交"的回测，
        在高换手策略上会系统性高估收益。
  \item \textbf{最小报价单位与整手}：低价股、高价股的实际可交易性与理论权重不同。
  \item \textbf{流动性摩擦}：Amihud 非流动性、买卖价差，直接影响可实现收益。
\end{itemize}

\subsection{信息扩散：你知道得比别人早，或者理解得比别人深}

财报公布后的漂移（PEAD）、分析师预期修正、公告文本的情绪变化，都属于这一类。
它的前提是\textbf{严格的时间点对齐}：你必须在信息真正公开的那一刻之后才能使用它。
用报告期（而不是公告日）对齐财务数据，是这一族因子里最常见的致命错误。

\subsection{拥挤与衰减：alpha 会自己消失}

任何被公开、被产品化、被资金追逐的规律，都会因为交易而衰减。这不是理论推测：

\begin{itemize}
  \item 小市值因子在 2017 年之前是 A 股最强的因子之一，2017 年之后随着监管与
        注册制扩容大幅削弱；
  \item 经典的中长期动量在 A 股长期偏弱甚至为负，而在美股相对稳定；
  \item 高频量价因子在 2021 年量化规模快速扩张后被显著压缩。
\end{itemize}

\begin{keybox}[判断一个因子能不能用的第一问]
它的收益来自\textbf{承担了什么风险}、\textbf{纠正了谁的错误}、还是\textbf{利用了哪条规则}？
如果三个都答不上来，那它大概率只是你在历史上挖到的一段噪声。
\end{keybox}

\section{机器学习到底贡献了什么}

这是全书最需要说清楚的一点。很多人以为机器学习能"发现人看不见的规律"，
但在日频横截面选股这个问题上，它的真实贡献要朴素得多：

\begin{enumerate}
  \item \textbf{非线性组合}。把 40 个因子线性加权，等于假设"每个因子的边际贡献恒定"。
        GBDT 能学到"低波因子在小市值股票里更有效"这类交互，从而提升一点 IC。
  \item \textbf{自动特征工程}。树模型能自动构造阈值与交叉项，省去人工试错。
  \item \textbf{对噪声的鲁棒性}。正则化与集成让模型不至于被单个极端样本带偏。
\end{enumerate}

它\textbf{不能}做的事同样清楚：

\begin{itemize}
  \item 它不能凭空创造新的经济来源。如果 40 个因子本身没有信息，模型也变不出来。
  \item 它不能对抗标签的噪声。日频收益的信噪比极低，模型很容易学到噪声。
  \item 它不能免除验证的责任。恰恰相反，模型越复杂，过拟合的自由度越大，
        需要的验证越严格。
\end{itemize}

\begin{center}
\begin{tabular}{@{}p{0.44\textwidth}p{0.46\textwidth}@{}}
\toprule
\textbf{机器学习真正擅长的} & \textbf{它做不到的} \\
\midrule
特征交互与非线性 & 无中生有的 alpha \\
多因子加权与降噪 & 替代严格的时间点对齐 \\
高频数据的模式识别 & 消除交易成本 \\
大规模横截面的批量预测 & 保证样本外有效 \\
\bottomrule
\end{tabular}
\end{center}

业界常引用的一个经验量级：在严格执行样本外验证之后，日频全 A 股的
模型 Rank IC 大致落在 $0.03$--$0.08$ 区间，年化 ICIR 在 $1$--$4$ 之间。
超过 $0.15$ 的日频 Rank IC，你首先应该怀疑的是数据泄漏，而不是庆祝。

\section{A 股的结构性特征}

把美股的那套方法直接搬到 A 股，会在四个地方翻车：

\subsection{T+1 与涨跌停}

T+1 意味着"当日买入不能当日卖出"，因此：
\begin{itemize}
  \item 标签不能用"当日收盘买、当日收盘卖"构造；
  \item 组合的换手约束更强，日频调仓的实际成本高于理论值。
\end{itemize}

涨跌停（主板 $\pm10\%$、创业板与科创板 $\pm20\%$、北交所 $\pm30\%$、ST $\pm5\%$）
意味着：
\begin{itemize}
  \item 涨停当天大概率买不进（尤其一字板），跌停当天卖不出；
  \item 如果不模拟这一点，反转类策略的回测收益会被系统性高估——
        因为"跌得最惨的股票"往往正是跌停卖不出的那些。
\end{itemize}

\subsection{散户结构与情绪波动}

A 股个人投资者成交占比长期偏高，导致：
\begin{itemize}
  \item 短期反转、换手率、情绪类因子比成熟市场更有效；
  \item 因子收益的波动更大，尾部风险更厚。
\end{itemize}

\subsection{注册制扩容与退市常态化}

2019 年科创板、2020 年创业板注册制、2021 年北交所，使得：
\begin{itemize}
  \item 股票数量从 2010 年的约 2000 只增长到 2026 年的约 5600 只；
  \item 小市值不再稀缺，"壳价值"逻辑失效；
  \item 退市股增多，\textbf{幸存者偏差}的杀伤力比过去更大。
\end{itemize}

\subsection{量化监管与交易约束}

2023 年之后对程序化交易的监管趋严，高频策略的容量被压缩。
这对个人研究者的启示是：\textbf{不要把研究建立在"极高换手"的假设上}，
那条路在 A 股的制度环境里越来越窄。

\section{哪些信号还在，哪些已经衰减}

把上文压缩成一张可执行的判断表。这是本书后续所有实证工作的先验：

\begin{center}
\begin{longtable}{@{}p{0.20\textwidth}p{0.13\textwidth}p{0.55\textwidth}@{}}
\toprule
\textbf{因子族} & \textbf{当前状态} & \textbf{判断依据} \\
\midrule
\endhead
短期反转（1--20 日） & 仍有效 & A 股最强的横截面效应，散户过度反应的直接结果 \\
低波动 / 特质波动 & 仍有效 & 杠杆约束 + 彩票偏好，跨市场稳健 \\
换手率 / 流动性 & 仍有效 & 拥挤度与关注度的代理变量 \\
价值（EP/BP） & 中等 & 大中盘较稳，小盘受"价值陷阱"干扰 \\
质量（ROE/现金流） & 中等 & 需要配合估值使用，单独使用收益有限 \\
规模（小市值） & 显著衰减 & 2017 年后壳价值消失，注册制进一步压缩 \\
中长期动量 & 弱甚至为负 & A 股中长期动量长期不稳定 \\
高频量价 & 衰减中 & 量化规模扩张 + 监管趋严 \\
\bottomrule
\end{longtable}
\end{center}

\begin{warnbox}[方向不是常识，是实证问题]
很多教程会写"动量因子：过去涨得多的股票未来继续涨"。在 A 股日频上，
这句话很可能是反的。本书第~\ref{ch:factor} 章的实测会显示：20 日动量的 IC 为负，
即\textbf{反转主导}。任何关于因子方向的断言，都必须在自己的样本上重新验证一次。
\end{warnbox}

\section{本章小结}

\begin{enumerate}
  \item 漂亮的回测曲线是极值统计量的自然产物，没有"尝试次数"信息的夏普比率不构成证据。
  \item 可预测性只有五类来源：风险补偿、行为偏差、微观结构、信息扩散、拥挤衰减。
        说不出来源的因子，大概率是噪声。
  \item 机器学习的贡献是"更好的组合与降噪"，不是"发现新规律"。
        它同时放大了过拟合的风险。
  \item A 股的 T+1、涨跌停、散户结构、注册制扩容，决定了哪些方法可行、哪些不可行。
  \item 因子方向必须实证，不能照搬教科书。
\end{enumerate}
